\chapter{Conclusiones}\label{cap:conclusiones}

\emph{Capítulo pendiente de desarrollo.} Las conclusiones definitivas de esta memoria dependen de los resultados de validación descritos como pendientes en la Sección~\ref{sec:validacion} (reevaluación heurística del prototipo y prueba de usabilidad con usuarios reales), que a su vez determinan el contenido del Capítulo~\ref{cap:discusion}. Redactar conclusiones sobre la efectividad del prototipo antes de contar con esa evidencia sería anticipar un resultado no verificado, lo que esta memoria busca evitar de forma consistente con el propio principio de transparencia que defiende para el SAE.

\section*{Lo que puede afirmarse hasta ahora}
\addcontentsline{toc}{section}{Lo que puede afirmarse hasta ahora}

Sin adelantar conclusiones sobre efectividad, el trabajo realizado hasta la fecha permite sostener lo siguiente, con el respaldo de los Capítulos~\ref{cap:introduccion} a~\ref{cap:resultados}:

\begin{enumerate}
    \item La brecha de transparencia del SAE real está documentada con evidencia propia (diagnóstico heurístico, Sección~\ref{sec:res-diagnostico}) y no es una percepción subjetiva: el sitio informativo obtuvo 0\,\% de cumplimiento en transparencia y apertura, y 0\,\% en inclusión.
    \item Existe un cuerpo consistente de literatura científica (Capítulo~\ref{cap:marco}) que identifica principios de diseño concretos y transferibles —divulgación progresiva, explicabilidad contextualizada, controles interactivos con retroalimentación— para abordar brechas de este tipo en sistemas algorítmicos de diversos dominios.
    \item Es técnicamente posible construir un prototipo funcional que implemente esos principios sin modificar el algoritmo de asignación subyacente: el prototipo desarrollado en este trabajo aborda, a nivel de código, la totalidad de los puntos aplicables de la matriz de verificación derivada del diagnóstico original y sus ampliaciones posteriores (Sección~\ref{sec:res-prototipo}). Se trata de una autoevaluación de implementación —cuyo conteo agregado quedó pendiente de re-tabular con criterio explícito, según la nota de integridad de la Sección~\ref{sec:res-prototipo}—, no de una medición de cumplimiento validada externamente.
    \item El sitio oficial del SAE, dieciséis meses después del diagnóstico original, ha cerrado de forma independiente solo dos de las seis brechas prioritarias identificadas (Sección~\ref{sec:res-sitio-oficial}), lo que sugiere que el problema de investigación de esta memoria sigue vigente y no ha sido resuelto por otra vía.
\end{enumerate}

\section*{Lo que falta para concluir el trabajo}
\addcontentsline{toc}{section}{Lo que falta para concluir el trabajo}

En orden de prioridad:

\begin{enumerate}
    \item \textbf{Ejecutar la reevaluación heurística del prototipo} con el mismo instrumento SISIB usado en el diagnóstico original, siguiendo el plan de cinco fases descrito en la Sección~\ref{sec:validacion}, e idealmente con triangulación de un segundo evaluador.
    \item \textbf{Ejecutar la prueba de usabilidad con usuarios reales} sobre el prototipo completo (no solo sobre el módulo de explicación del resultado, a diferencia del estudio previo mencionado en la Sección~\ref{sec:fase2}), replicando un esquema de tareas dirigidas y cuestionario de percepción comparable.
    \item \textbf{Redactar el Capítulo~\ref{cap:discusion} (Discusión)} interpretando esos resultados a la luz de la literatura del Capítulo~\ref{cap:marco} y del estado del sitio oficial reportado en la Sección~\ref{sec:res-sitio-oficial}.
    \item \textbf{Redactar este capítulo de Conclusiones} en su versión definitiva, incluyendo las contribuciones reales del trabajo, sus limitaciones confirmadas empíricamente, y las líneas de trabajo futuro —entre ellas, la colaboración con el Ministerio de Educación y la ampliación a otros perfiles de usuario— que ya se perfilan en la Sección~\ref{sec:limitaciones} pero que solo pueden fijarse con propiedad una vez conocidos los resultados.
\end{enumerate}
